\section{Rome, Family, and the Tribunal of Milan}

Ambrose belonged by birth to the governing aristocracy of the Christianizing Roman Empire. His father, also named Ambrosius, held the office of praetorian prefect of the Gauls, a civilian jurisdiction stretching across Gaul, Spain, and Britain, with its chief administrative seat at Trier (Augusta Treverorum) on the Moselle. Ambrose was the youngest of three children, born after his elder sister Marcellina and his elder brother Satyrus. No surviving letter or treatise of Ambrose names his exact birthplace or birth year; Trier is the probable birthplace because it was the prefect's seat, and c.~339--340 is the widely received birth range. An older alternative places his birth around 333, turning on whether an autobiographical allusion to age amid ``barbaric commotions and the storms of war'' in \work{Epistle} 59 refers to the invasion of Maximus in 387 or of Eugenius in 393--394; the uncertainty remains unresolved in the primary record.

His family was Christian before his birth and counted both state officials and confessors or martyrs among its kin. After his father died---probably while Ambrose was still a boy---his widowed mother brought her three children back to Rome, no earlier than Christmas 353 (and more probably 354). One domestic scene from those Roman years survives in Ambrose's own hand. Addressing Marcellina more than twenty years later in his treatise \work{Concerning Virgins}, he recalled the Christmas Day when Pope Liberius (bishop of Rome, 352--366) conferred the veil of consecrated virginity upon her in the basilica of St.~Peter before a throng of virgins. Ambrose preserved the substance of Liberius's address as Marcellina had long talked it over with him:

\begin{quote}
``You,'' said he, ``my daughter, have desired a good espousal. You see how great a crowd has come together for the birthday of your Spouse, and none has gone away without food. This is He, Who, when invited to the marriage feast, changed water into wine. He, too, will confer the pure sacrament of virginity on you who before were subject to the vile elements of material nature. This is He Who fed four thousand in the wilderness with five loaves and two fishes.'' \ldots\ And now He has called many to your espousal, but it is not now barley bread, but the Body from heaven which is supplied. (\work{De Virginibus} III.1.1)
\end{quote}

Ambrose placed strong Nicene Christology directly on Liberius's lips---``God of God, born of a Virgin, \ldots\ Light of Light, not unequal to His Father'' (\work{De Virginibus} III.1.2)---a striking memorial given that Liberius would soon suffer exile under Constantius~II in 355 and later draw bitter criticism from Hilary, Athanasius, and Jerome for temporizing under imperial pressure. Within Ambrose's household, however, Liberius remained ``of blessed memory,'' and Marcellina's silent consecrated life at Rome formed the steady spiritual anchor of her brothers' world.

At Rome, Ambrose and Satyrus received the standard upper-class education in Latin grammar, classical rhetoric, Roman law, and Greek---a mastery of Greek unusually fluent for a fourth-century Westerner, which would later let Ambrose read Philo, Origen, Athanasius, Didymus the Blind, and Basil of Caesarea directly in the original. Though raised in a devout Christian house, neither brother received baptism in youth, following a widespread fourth-century custom among public men who postponed the sacrament out of dread of post-baptismal sin while holding office.

Both brothers entered the bar and rose through the imperial administration attached to the praetorian prefecture of Italy under Sextus Claudius Petronius Probus. Ambrose's own funeral oration for Satyrus recalls his brother's ``well-known eloquence in his forensic duties'' in ``the hall of justice of the high prefecture'' (\work{De Excessu Fratris Satyri} I.49) and remarks that both brothers accepted secular honors without craving them, ``not because the acquisition of them was to be desired, but that there might be no appearance of paltry dissimulation'' (\work{De Excessu} I.25). Around 370, Probus appointed Ambrose governor of the northern Italian province of Aemilia and Liguria with the senatorial rank of \latin{consularis}---one of thirty-seven consular governorships across the Empire---seated at Milan.

\witnesslimit{The famous parting charge attributed to Prefect Probus---``Go and act, not as a judge, but as a bishop''---derives from Paulinus of Milan's \work{Vita Ambrosii} (c.~422; summarized in H.~de Romestin's Prolegomena). No contemporary writing of Ambrose repeats the remark, though Ambrose's own portrait of his brother's provincial magistracy as ``rather their parent than a judge'' (\work{De Excessu} I.58) reflects the same late-Roman ideal of equitable rule.}

As \latin{consularis}, Ambrose presided over civil and criminal courts, tax assessments, municipal order, and public security in a province whose capital housed the western imperial court. His training taught him the structure of imperial rescripts, the etiquette of senatorial petitions, the vocabulary of Roman law, and the habit of public command. None of his colleagues expected that those civic gifts were preparing him not for a higher prefecture, but for a bishop's chair.
